\documentclass{book}
\usepackage{amssymb,amsmath}
\usepackage{polyglossia}
\setmainlanguage{spanish} % Idioma principal
\usepackage{theorem}
\usepackage{times}
\usepackage{array}
\usepackage{graphicx}
\usepackage{hyperref}
\usepackage{multirow}
\usepackage{fancyhdr}
%\usepackage[cp1252]{inputenc}
\usepackage{hhline}
\usepackage{multicol}
\usepackage[a4paper,driver=xetex,top=4.5cm,head=4.5cm, bottom=2cm,%
layouthoffset=0mm, left=2.5cm, right=2.5cm,marginparwidth=0cm]{geometry}
%\usepackage{bm}
%\usepackage{tabular}
\usepackage{fontspec}
 \usepackage[breakable,many]{tcolorbox}
\defaultfontfeatures{Ligatures=TeX}
\usepackage{empheq}
 \setromanfont{Roboto Condensed}
 \usepackage{float}
 \usepackage{mathrsfs} 
%
%
% \renewcommand{\familydefault}{\sfdefault}
%\renewcommand{\familydefault}{\sfdefault}

%%%%%%%%%Estilo de la pagina%%%%%%%%%%%%%%%%%%%%%%%%%%%%%%%%%%%
%%%%%%%%%%%%%%%%%%%%%%%%%%%%%%%%%%%%%%%%%%%%%%%%%%%%%%%%%%%%%%%%%%
% \newcounter{ejer}
% 
% {\theorembodyfont{\normalfont}
% \newtheorem{ejercicio}[ejer]{Ejercicio}}

\newcommand{\rr}{\mathbb{R}}
\newcommand{\qq}{\mathbb{Q}}
\newcommand{\nn}{\mathbb{N}}



\DeclareMathOperator{\atan2}{atan2}
%\DeclareMathOperator{\sen}{sen}
\DeclareMathOperator{\sign}{sign}
\DeclareMathOperator{\sn}{sn}
\DeclareMathOperator{\SO}{SO}
%\DeclareMathOperator{\arcsen}{arcsen}
\DeclareMathOperator{\Or}{O}

\usepackage[framemethod=TikZ]{mdframed}
%%%%%%%%%%%%%%%%%%%%%%%%%%%%%%
%Theorem

%% Ejercicio
\newcounter{ejer} \setcounter{ejer}{0}
\renewcommand{\theejer}{\arabic{ejer}}
\newenvironment{ejer}[2][]{%
\vspace{5pt}
\refstepcounter{ejer}%
\ifstrempty{#1}%
{%
% \mdfsetup{%
% frametitle={%
% \tikz[baseline=(current bounding box.east),outer sep=-0pt]
% \node[anchor=east,rectangle,fill=green!50]
{\noindent\bfseries Ejercicio~\theejer}.}
%
{%
% \mdfsetup{%
% frametitle={%
% \tikz[baseline=(current bounding box.east),outer sep=0pt]
% \node[anchor=east,rectangle,fill=green!50]
{\noindent\bfseries  Ejercicio~\theejer:~#1};}%
%
%\mdfsetup{innertopmargin=10pt,linecolor=green!50,%
%linewidth=2pt,topline=true,%
%frametitleaboveskip=\dimexpr-\ht\strutbox\relax
%}
%\begin{mdframed}[]
\relax%
\label{#2}}{\vspace{5pt}}%\end{mdframed}}

%Theorem
\newcounter{theo}[chapter] \setcounter{theo}{0}
\renewcommand{\thetheo}{\arabic{section}.\arabic{theo}}
\newenvironment{theo}[2][]{%
\refstepcounter{theo}%
\ifstrempty{#1}%
{\mdfsetup{%
frametitle={%
\tikz[baseline=(current bounding box.east),outer sep=0pt]
\node[anchor=east,rectangle,fill=blue!20]
{\strut Teorema~\thetheo};}}
}%
{\mdfsetup{%
frametitle={%
\tikz[baseline=(current bounding box.east),outer sep=0pt]
\node[anchor=east,rectangle,fill=blue!20]
{\strut Teorema~\thetheo:~#1};}}%
}%
\mdfsetup{innertopmargin=10pt,linecolor=blue!20,%
linewidth=2pt,topline=true,%
frametitleaboveskip=\dimexpr-\ht\strutbox\relax
}
\begin{mdframed}[]\relax%
\label{#2}}{\end{mdframed}}
%%%%%%%%%%%%%%%%%%%%%%%%%%%%%%
%Lemma
\newcounter{lem}[chapter] \setcounter{lem}{0}
\renewcommand{\thelem}{\arabic{section}.\arabic{lem}}
\newenvironment{lem}[2][]{%
\refstepcounter{lem}%
\ifstrempty{#1}%
{\mdfsetup{%
frametitle={%
\tikz[baseline=(current bounding box.east),outer sep=0pt]
\node[anchor=east,rectangle,fill=green!20]
{\strut Lemma~\thelem};}}
}%
{\mdfsetup{%
frametitle={%
\tikz[baseline=(current bounding box.east),outer sep=0pt]
\node[anchor=east,rectangle,fill=green!20]
{\strut Lemma~\thelem:~#1};}}%
}%
\mdfsetup{innertopmargin=10pt,linecolor=green!20,%
linewidth=2pt,topline=true,%
frametitleaboveskip=\dimexpr-\ht\strutbox\relax
}
\begin{mdframed}[]\relax%
\label{#2}}{\end{mdframed}}
%%%%%%%%%%%%%%%%%%%%%%%%%%%%%%
%% Definicion
\newcounter{defini}[chapter] \setcounter{defini}{1}
\renewcommand{\thedefini}{\arabic{section}.\arabic{defini}}
\newenvironment{definicion}[2][]{%
\refstepcounter{defini}%
\ifstrempty{#1}%
{\mdfsetup{%
frametitle={%
\tikz[baseline=(current bounding box.east),outer sep=0pt]
\node[anchor=east,rectangle,fill=green!20]
{\strut Definición~\thedefini};}}
}%
{\mdfsetup{%
frametitle={%
\tikz[baseline=(current bounding box.east),outer sep=0pt]
\node[anchor=east,rectangle,fill=green!20]
{\strut Definición~\thedefini:~#1};}}%
}%
\mdfsetup{innertopmargin=10pt,linecolor=green!20,%
linewidth=2pt,topline=true,%
frametitleaboveskip=\dimexpr-\ht\strutbox\relax
}
\begin{mdframed}[]\relax%
\label{#2}}{\end{mdframed}}

%Proof
\newenvironment{prf}{\noindent\emph{Dem.}}{$\square$ \newline\vspace{5pt}}


%Corolario
\newcounter{cor}[chapter] \setcounter{cor}{0}
\renewcommand{\thecor}{\arabic{section}.\arabic{cor}}
\newenvironment{cor}[2][]{%
\refstepcounter{cor}%
\ifstrempty{#1}%
{\mdfsetup{%
frametitle={%
\tikz[baseline=(current bounding box.east),outer sep=0pt]
\node[anchor=east,rectangle,fill=green!20]
{\strut Corolario~\thelem};}}
}%
{\mdfsetup{%
frametitle={%
\tikz[baseline=(current bounding box.east),outer sep=0pt]
\node[anchor=east,rectangle,fill=green!20]
{\strut Corolario~\thelem:~#1};}}%
}%
\mdfsetup{innertopmargin=10pt,linecolor=green!20,%
linewidth=2pt,topline=true,%
frametitleaboveskip=\dimexpr-\ht\strutbox\relax
}
\begin{mdframed}[]\relax%
\label{#2}}{\end{mdframed}}

\tcbset{highlight math style={enhanced,
  colframe=red!60!black,colback=yellow!50!white,arc=4pt,boxrule=1pt,
  drop fuzzy shadow}}
  
  
  
  
  
  
  


\pagestyle{fancyplain}

 \renewcommand{\sectionmark}[1]
                 {\markright{\thesection\ #1}}


% \lhead[\fancyplain{}{\bfseries\thepage}]
%       {\fancyplain{}{\bfseries\rightmark}}
%
 \rhead[\fancyplain{}{\bfseries\leftmark}]{\fancyplain{}{\bfseries}}




 \lhead[\fancyplain{}{ \includegraphics[scale=.3]{EscudoUNLPam.png}}]{\fancyplain{}{ \includegraphics[scale=.3]{EscudoUNLPam.png}}}

\cfoot{}





  
  
  
  
  
  
  
\begin{document}


\hyphenation{excen-tri-ci-dad}


\begin{large}
\begin{bfseries} % \begin{scshape}
        \noindent Depto de Matem\'atica.\\
        Primer Cuatrimestre de 2022\\                                                                                                                                                                                                                                                                                                                                                
        Teoría de la Medida \\
        Práctica 5: Integral de Lebesgue

%\end{scshape}
\end{bfseries}
\end{large}
\par\noindent\rule{\textwidth}{.5pt}

%\begin{color}{red}
\begin{ejer}{}  Por medio de la definición de integral de función no negativa, hallar el valor de las integrales
\[
 \int_0^\infty \sum_{n=1}^\infty \chi_{\left[n,n+\frac{1}{2^n}\right)}\quad\text{y}\quad \int_0^\infty \sum_{n=1}^\infty \alpha_n\chi_{\left[n,n+\frac{1}{2}\right)},
\]
donde $\alpha_n\geq 0$. 
 
\end{ejer}
%\end{color}

%\begin{color}{red}
\begin{ejer}[Criterio de Lebesgue para la integrabilidad Riemann]{} El criterio establece que si $f:[a,b]\to\rr$ es acotada entonces  son equivalentes
\begin{enumerate}
\item[a)] $f$ es integrable Riemann.
 \item[b)] Si $Z$ es el conjunto de puntos donde $f$ es discontinua entonces $m_*(Z)=0$. 
\end{enumerate}
\emph{Ayuda:} Para a)$\Rightarrow$ b) observar que 
\[
 Z=\{x:\omega(f;x)>0\}=\bigcup_{n=1}^{\infty}S_n:= \bigcup_{n=1}^{\infty}\left\{x:\omega(f;x)\geq \frac{1}{n}\right\}.
\]
Luego usar que $m_*(S_n)\leq c_e(S_n)$.

Para b)$\Rightarrow$ a). Notar que $m_*(S_n)=0$. Luego, para todo $\varepsilon>0$ se puede hallar una colección de intervalos abiertos y disjuntos $I_j$ tales que $S_n\subset \bigcup\limits_{j=1}^\infty I_j$ y   $\sum\limits_j m(I_j)<\varepsilon$.  Demostrar que $S_n$ es compacto y de allí se cubre por finitos $I_j$. Finalmente invocar la caracterización de Hankel.
\end{ejer}
%\end{color}




\begin{ejer}{}
  %Mostrar que la función $x^{p-1} e^{-x}$ es %integrable sobre $(0,\infty)$ si y sólo si $p>0$. %\textcolor{red}{Quizás es más importante que %aprenda cuando las funciones $1/x^p$ son %integrales en $(0,1)$ y $(1,\infty)$. ¿no?}
  Analizar para qu\'e valores de $p$ las funciones $\frac{1}{x^p}$ son integrables en $(0,1)$ y $(1,\infty)$.
	\end{ejer}
  
	\begin{ejer}{} 
	Probar que la función $\frac{\text{sen}\,x}{x}$ no es integrable sobre $(0,\infty)$, aunque existe el límite
\linebreak  $\lim\limits_{R\rightarrow \infty}\int_0^R \frac{\text{sen}\,x}{x}\,dx.$
\end{ejer}
	
%	\begin{ejer}{}
%	Probar que la integral $\int_0^1 \frac{1}{x}\text{sen} \left(\frac{1}{x}\right)\,dx$ 
	%existe como integral impropia de Riemann %pero no existe como integral de Lebesgue.
	%\end{ejer}

%e un intervalo  infinito (tal integral
	%s\'olo puede existir en el sentido impropio). 
	%Demostrar que $f$ es integrable de Lebesgue sobre el mismo intervalo si y s\'olo si la integral
	%impropia converge absolutamente.
 % \end{ejer}
	
	
	\begin{ejer}{}
	Probar, usando el Teorema de la Convergencia Mayorada, la fórmula
  $$\lim\limits_{n\rightarrow \infty} \int_0^1 \frac{n^{\frac{3}{2}} x}{1+n^2 x^2}\,dx=0. $$
	\end{ejer}


  \begin{ejer}{}
	Usando integración término a término probar que 
  $$\int_0^{\infty} \frac{x}{e^x-1}\,dx=\sum_{n=1}^{\infty}\frac{1}{n^2}. $$
  \end{ejer}
	
	\begin{ejer}{}
	\begin{enumerate}
	\item Sup\'ongase que $f\geq 0$, integrable de Riemann en $[0,a]$ $\forall a>0$ y que adem\'as tiene integral impropia. 
	Probar que $f$ es integrable de Lebesgue en $[0,+\infty)$.
	%\textcolor{red}{¿Capaz que sea en $[0,+\infty)$, no?}.
	\item Sup\'ongase que $f\geq 0$, integrable de Riemann en $[a+\epsilon,b]$ y que existe la integral impropia.
	Probar que $f$ es integrable de Lebesgue en $[a,b]$.
	\end{enumerate}
	\end{ejer}
	
  \begin{ejer}{} 
  Sea $f$ una función medible no negativa sobre $\rr$ y sea $(E_k)$ una sucesión creciente de conjuntos
  cuya unión es $E$.
	\begin{enumerate}
  \item Probar que $\int_{E} f=\lim\limits_{k \rightarrow \infty} \int_{E_k} f$.
  \item Extender a cualquier función $f$ integrable sobre $E$.
	\end{enumerate}
	\end{ejer}

  

  \begin{ejer}{}
	Si se considera la sucesión de funciones $f_n(x)=n \chi_n(x)$ en el intervalo $0\leq x\leq 1$,
  donde $\chi_n$ es la función característica del intervalo $\left(0,\frac{1}{n}\right)$, ?`es posible que exista
  una función $g(x)$ integrable en dicho intervalo, tal que $f_n(x) \leq g(x)$ para cualquier $n$ y cualquier
  $x$?
	\end{ejer}


 %\begin{ejer}{} \textcolor{red}{Yo creo Sonia que %esto estaría bueno dejarlo para los $L^p$}
%	Si $\varphi(x) f(x)$ es integrable sobre $E$ para %cualquier función $f$ integrable sobre $E$, entonces
%  existe una constante finita $C$, tal que $|\varphi(x)|\leq C$ en c.t.p $x$ de $E$.
%\end{ejer}
  
  \begin{ejer}{} 
	Sea $f$ una función medible no negativa sobre $\rr^1$ tal que $\int_a^b f(x)\,dx>0$ siempre
  que $a<b$, ?`puede concluirse que $f(x)>0$ en c.t.p $x$?
	\end{ejer}
 
%\begin{ejer}{}\textcolor{red}{Capaz lo trabaje en el teórico}
%	\textbf{Absoluta continuidad de la integral de Lebesgue.}
%	Sup\'ongase que $f$ es integrable de Lebesgue sobre $E$. 
%	Probar que para cada $\epsilon>0$ existe %$\delta>0$ tal que $\int_E |f|\,dx<\epsilon$ siempre %que $m(E)<\delta$.
	%\end{ejer}
	
  \begin{ejer}{} 
  Probar que el Teorema de la Convergencia Mayorada se extiende a una familia de funciones medibles $f_t(x)$,\;  
  $a<t<b$, dependiente de un parámetro real $t$, de la manera siguiente:
   \\
   Supongamos que $\tau \in (a,b)$ y que en cada punto de $E$ existe el límite 
   $f(x)=\lim\limits_{t \rightarrow \tau} f_t(x)$.
   Si existe una función $\Phi(x)$ integrable sobre $E$, tal que $|f_t(x)|\leq \Phi(x)$ para $x \in E$ y $a<t<b$;
   entonces $f$ es integrable sobre $E$ y además 
   $$\int_E f(x)\,dx=\lim_{t \rightarrow \tau}\int_E f_t(x)\,dx. $$
	\end{ejer}
	
	\begin{ejer} [Derivación de una integral paramétrica]{}*
   Supongamos que la integral 
   $$ \varphi(t)=\int_E f(t,x)\,dx\;\;\;\;\;\;\;(a<t<b),$$ 
   existe para cada $t\in (a,b)$; que $f(t,x)$ es derivable con respecto
   a $t$ y existe una función $g(x)$ integrable sobre $E$, tal que 
   $$\Bigl|\frac{\partial f(t,x)}{\partial t}\Bigr|\leq g(x)$$
   para $x \in E$ y $a<t<b$.
   Probar que $\varphi$ es derivable y además 
   $$\varphi'(t)=\int_E \frac{\partial f(t,x)}{\partial t}\,dx $$ para $a<t<b$.
	\\
   \underline{Sugerencia:} Escribir el cociente $\frac{\varphi(t+h)-\varphi(t)}{h}$, emplear el Teorema del Valor
   Medio del Cálculo Diferencial y el ejercicio anterior.
	\end{ejer}

\begin{ejer}[Transformada de Fourier]{}*
 Si $f$ es integrable sobre $\rr$, la función 
  $$g(t)=\int_{-\infty}^{\infty} e^{i t x} f(x)\,dx $$
   es acotada y uniformemente continua.
   \\
   Si $x^k.f(x)$ es integrable, entonces $g$ es de clase $C^k$ y además 
   $$g^{(k)}(t)=\int_{-\infty}^{\infty} e^{i t x} (i x)^k  f(x)\,dx.$$
 \end{ejer}
   
   
 
  
%\begin{ejer}{} \textcolor{red}{Como vamos a adoptar el enfoque de Fava y Zó no sería adecuado este ejercicio}
 %Si $(\rr,\mathscr{B})$ denota el espacio medible que consiste de los números reales junto con los conjuntos 
%de Borel, mostrar que todo subconjunto abierto de $\rr\times \rr$ pertenece a $\mathscr{B}\times \mathscr{B}$.
%De hecho, esta $\sigma$-álgebra es la sigma-álgebra generada por los subconjuntos abiertos de $\rr\times \rr$.
%(En otras palabras, $\mathscr{B}\times \mathscr{B}$ es el álgebra de Borel de $\rr \times \rr$).
%\end{ejer}


\begin{ejer}{} 
Integrar la función no negativa $xe^{-x^2(1+y^2)}$ sobre el conjunto $(0,\infty)\times (0,\infty)$
en dos formas diferentes. Concluir a partir de sus cálculos que 
\[
\int_0^{\infty} e^{-t^2}\,dt=\frac{\sqrt{\pi}}{2}.
\]
\end{ejer}


\begin{ejer}{} 
\begin{enumerate}
\item Integrar la función $\text{sen} x e^{-xy}$ sobre el conjunto $(0,a)\times (0,\infty)$.
Mostrar que 
\begin{equation}\label{eq:fubini-senx/x}
\int_0^a \frac{\text{sen} x}{x}\,dx=
\frac{\pi}{2}-\cos a\int_0^{-\infty} \frac{e^{-ay}}{1+y^2}\,dy
-\sen a \int_0^{\infty} \frac{y e^{-ay}}{1+y^2}\,dy
\end{equation}
\item Mostrar que 
\[
\int_0^{\infty }\frac{\text{sen} x}{x}\,dx=\frac{\pi}{2}
\]
\underline{Ayuda:} 
Demostrar que las integrales de miembro derecho de \eqref{eq:fubini-senx/x} tienden a 0 cuando $a \to \infty$.
\end{enumerate} 
\end{ejer}


\begin{ejer}{}
\begin{enumerate}
\item Mostrar que
\[
\int_0^{\infty} \frac{\text{sen} a x}{x}\,dx= 
\left
\{\begin{array}{rll}
\frac{\pi}{2}&si&a>0
\\
0&si&a=0
\\
-\frac{\pi}{2}&si&a<0
\end{array}
\right.
\] 
\item Integrando por partes en el item anterior mostrar que
\[
\int_0^{\infty}\frac{1-\cos ax}{x^2}\,dx=\frac{\pi}{2}|a|.
\]
\item Probar que $\frac{1-\cos ax}{x^2}$ es una función integrable en $(0,\infty)$.
\end{enumerate}
\end{ejer}


\begin{ejer}{}
 Por integración de $e^{-xy}$ sobre una región apropiada mostrar que 
\[
\int_0^{\infty}\frac{e^{-ax}-e^{-bx}}{x}\,dx=\ln\frac{b}{a}\;\; \mbox{  si   }a>0, \;\;b>0.
\]
\end{ejer}


%\textcolor{red}{Los siguientes ejericios me gustan todos, pero ¿No quedará demasiado larga la práctica?. ¿NO te parece que asterisquemos 23,25,26,27?}

\begin{ejer}{}*
 Sean $a_i>0$ para $1\leq i \leq n$ y $J=(0,1)\times \ldots \times (0,1)$.
Probar que 
\[
\int_J  \frac{1}{x_1^{a_1}+x_2^{a_2}+\ldots+x_n^{a_n}}\,dx<\infty 
\Longleftrightarrow
\sum\limits_{i=1}^{n}\frac{1}{a_i}>1.
\]
\\
\underline{Ayuda:}
Sea $G_i=\{x \in J | x_j^{a_j} \leq x_i^{a_i} \mbox{ para todo  } j \}$. \\
Notar que para $x \in G_i$, $x_i^{a_i}\leq x_1^{a_1}+x_2^{a_2}+\ldots+x_n^{a_n}\leq n x_i^{a_i}$
y $J=\bigcup\limits_{i=1}^{n} G_i$. \\
Calcular 
\[
\int_G \frac{dx}{x_i^{a_i}}=\int_0^1 x_i ^{a_i\left(\sum\limits_{j=1}^n \frac{1}{a_j}-1\right)-1}\,dx_i,
\]
y usar el hecho de que $\int_0^1 t^{s-1}\,dt<\infty \Leftrightarrow s>0$.
\end{ejer}


\begin{ejer}{}
Sobre $(0,1)\times (0,1)$ considerar
\[
f(x,y)=\left\{
\begin{array}{rll}
x^{-2}&si&y<x<1
\\
-y^{-2}&si&x<y<1
\end{array}
\right.
\]
y probar que $\int_0^1\,dx\int_0^1 f(x,y)\,dy=1$ y $\int_0^1\,dy \int_0^1 f(x,y)\,dx=-1$.
\end{ejer}

\begin{ejer}{}*
 Mostrar que si $f(x,y)=\frac{(x^2-y^2)}{(x^2+y^2)^2}$, entonces
     $$\frac{\pi}{4}=\int_0^1 \,dx\int_0^1 f(x,y)\,dy\not = \int_0^1\,dy \int_0^1 f(x,y)\,dx=-\frac{\pi}{4}. $$
\end{ejer}

 \begin{ejer}{}* Mostrar que la función $f(x,y)=\frac{xy}{(x^2+y^2)^2}$ verifica
   $$\int_{-1}^1 \,dx\int_{-1}^1 f(x,y)\,dy=\int_{-1}^1\,dy \int_{-1}^1 f(x,y)\,dx $$
   a pesar de que $f$ no es integrable sobre el cuadrado $Q=[-1,1]\times [-1,1]$.
	\end{ejer}

 
\begin{ejer}{}* Una función no negativa $f$, definida sobre $\rr^1$, se llama una {\bf{densidad}}
   (de probabilidad) si su integral sobre toda la recta es igual a uno. Probar que: 
	\begin{enumerate}
	\item La convolución de dos densidades es otra densidad.
	\item Para cada $p>0$, la función 
     $$f_p(x)=\left\{ 
      \begin{array}{cc}
      \frac{1}{\Gamma(p)} x^{p-1} e^{-x} & \mbox{si $x>0$}\\
      0 & \mbox{si $x \leq 0$}
      \end{array}
      \right. $$
      es una densidad.
	 \item Si $p>0$ y $q>0$, entonces $f_p*f_q=f_{p+q}$ y además
      $$\int_0^1 (1-t)^{p-1} t^{q-1}\,dt=\frac{\Gamma(p) \Gamma(q)}{\Gamma(p+q)}. $$
      La integral biparamétrica del primer miembro se denota por $B(p,q)$ y se llama 
      {\bf{función beta}}.
	\end{enumerate}
	
\end{ejer}




	
	
	
	


%¿¿¿¿¿¿¿¿¿¿¿¿¿¿¿¿¿¿¿¿¿¿¿¿¿¿¿¿¿¿¿¿¿¿¿¿¿¿¿


\end{document}
